% Part: first-order-logic
% Chapter: sequent-calculus
% Section: derivations

\documentclass[../../../include/open-logic-section]{subfiles}

\begin{document}

\iftag{FOL}
      {\olfileid{fol}{seq}{der}}
      {\olfileid{pl}{seq}{der}}

\olsection{\usetoken{P}{derivation}}

\begin{explain}
આરંભિક સિક્વન્ટ કેવું હોય છે તે આપણે જણાવ્યું છે અને અનુમાનના નિયમો પણ
આપ્યા છે. સિક્વન્ટ કલનમાં !!^{derivation}s આ બંનેમાંથી આગમનાત્મક રીતે
ઉત્પન્ન થાય છે: દરેક !!{derivation} કાં તો એકલું આરંભિક સિક્વન્ટ હોય છે,
અથવા એક કે બે !!{derivation}s પછી આવેલું અનુમાન ધરાવે છે.
\end{explain}

\begin{defn}[$\Log{LK}$-!!{derivation}]
સિક્વન્ટ~$S$ની \emph{$\Log{LK}$-!!{derivation}} એ નીચેની શરતો
સંતોષતું સિક્વન્ટોનું સાન્ત વૃક્ષ છે:
\begin{enumerate}
\item વૃક્ષના સૌથી ઉપરના સિક્વન્ટો આરંભિક સિક્વન્ટો છે.
\item વૃક્ષનો સૌથી નીચેનો સિક્વન્ટ~$S$ છે.
\item વૃક્ષમાં $S$ સિવાયનો દરેક સિક્વન્ટ અનુમાનના કોઈ નિયમના યોગ્ય
  ઉપયોગનો આધાર-સિક્વન્ટ છે, અને તે નિયમનું ફલિત-સિક્વન્ટ વૃક્ષમાં તેની
  સીધી નીચે ઊભું છે.
\end{enumerate}
પછી $S$ને !!{derivation}નો \emph{અંતિમ સિક્વન્ટ} કહીએ છીએ અને કહીએ છીએ
કે $S$ \emph{$\Log{LK}$માં !!{derivable} છે} (અથવા
$\Log{LK}$-!!{derivable} છે).
\end{defn}

\begin{ex}
દરેક આરંભિક સિક્વન્ટ, દાખલા તરીકે $!C \Sequent !C$, એ !!a{derivation}
છે. તેના પર, માનો કે, $\LeftR{\Weakening}$ નિયમ લગાડીને આપણે નવી
!!{derivation} મેળવી શકીએ છીએ:
\begin{prooftree}
\Axiom$ \Gamma \fCenter \Delta $
\RightLabel{\LeftR{\Weakening}}
\UnaryInf$ !A, \Gamma \fCenter \Delta$
\end{prooftree}
જોકે નિયમ સામાન્ય અર્થમાં વાપરવાનો છે: નિયમમાંના $!A$ને કોઈપણ
!!{sentence}થી, દાખલા તરીકે~$!D$થી પણ, બદલી શકીએ છીએ. જો આધાર-સિક્વન્ટ
આપણા આરંભિક સિક્વન્ટ $!C \Sequent !C$ સાથે મળતું હોય, તો $\Gamma$ અને
$\Delta$ બંને માત્ર~$!C$ છે અને ફલિત-સિક્વન્ટ $!D, !C \Sequent !C$
થશે. તેથી નીચેનું !!a{derivation} છે:
\begin{prooftree}
\Axiom$ !C \fCenter !C $
\RightLabel{\LeftR{\Weakening}}
\UnaryInf$ !D, !C \fCenter !C$
\end{prooftree}
હવે આપણે બીજો કોઈ નિયમ, માનો કે $\LeftR{\Exchange}$, લગાડી શકીએ છીએ;
તે ડાબી બાજુનાં બે !!{sentence}sની અદલાબદલી કરવા દે છે. તેથી નીચેનું પણ
યોગ્ય !!{derivation} છે:
\begin{prooftree}
\Axiom$ !C \fCenter !C $
\RightLabel{\LeftR{\Weakening}}
\UnaryInf$ !D, !C \fCenter !C$
\RightLabel{\LeftR{\Exchange}}
\UnaryInf$ !C, !D \fCenter !C$
\end{prooftree}
નિયમ આ પ્રમાણે આપ્યો હતો:
\begin{prooftree}
\Axiom$ \Gamma, !A, !B, \Pi \fCenter \Delta $
\RightLabel{\LeftR{\Exchange}}
\UnaryInf$ \Gamma, !B, !A, \Pi \fCenter \Delta,$
\end{prooftree}
તેના ઉપરના ઉપયોગમાં $\Gamma$ અને $\Pi$ બંને ખાલી હતાં, $\Delta$ એ
$!C$ હતું, અને $!A$ તથા $!B$ની ભૂમિકા અનુક્રમે $!D$ અને~$!C$એ ભજવી
હતી. લગભગ એ જ રીતે આપણે એ પણ જોઈ શકીએ છીએ કે
\begin{prooftree}
\Axiom$ !D \fCenter !D $
\RightLabel{\LeftR{\Weakening}}
\UnaryInf$ !C, !D \fCenter !D$
\end{prooftree}
એ !!a{derivation} છે. હવે આ બંને !!{derivation}s લઈને
$\RightR{\land}$ વડે તેમને જોડી શકીએ છીએ. તે નિયમ આ હતો:
\begin{prooftree}
\Axiom$\Gamma \fCenter \Delta, !A$
\Axiom$ \Gamma \fCenter \Delta, !B$
\RightLabel{\RightR{\land}}
\BinaryInf$ \Gamma \fCenter \Delta, !A \land !B $
\end{prooftree}
આપણા કિસ્સામાં આધાર-સિક્વન્ટો તેમનાથી પૂરાં થતાં !!{derivation}sના
છેલ્લા સિક્વન્ટો સાથે મળવા જોઈએ. એટલે $\Gamma$ એ $!C, !D$ છે,
$\Delta$ ખાલી છે, $!A$ એ $!C$ છે અને $!B$ એ $!D$ છે. તેથી અનુમાન
યોગ્ય હોય તો તેનું ફલિત-સિક્વન્ટ $!C, !D \Sequent !C \land !D$ છે.
\begin{prooftree}
\Axiom$ !C \fCenter !C $
\RightLabel{\LeftR{\Weakening}}
\UnaryInf$ !D, !C \fCenter !C$
\RightLabel{\LeftR{\Exchange}}
\UnaryInf$ !C, !D \fCenter !C$
\Axiom$ !D \fCenter !D $
\RightLabel{\LeftR{\Weakening}}
\UnaryInf$ !C, !D \fCenter !D$
\RightLabel{\RightR{\land}}
\BinaryInf$ !C, !D \fCenter !C \land !D $
\end{prooftree}
અલબત્ત, આધાર-સિક્વન્ટોનો ક્રમ ઊલટાવી પણ શકીએ; ત્યારે $!A$ એ $!D$
અને $!B$ એ~$!C$ થશે.
\begin{prooftree}
\Axiom$ !D \fCenter !D $
\RightLabel{\LeftR{\Weakening}}
\UnaryInf$ !C, !D \fCenter !D$
\Axiom$ !C \fCenter !C $
\RightLabel{\LeftR{\Weakening}}
\UnaryInf$ !D, !C \fCenter !C$
\RightLabel{\LeftR{\Exchange}}
\UnaryInf$ !C, !D \fCenter !C$
\RightLabel{\RightR{\land}}
\BinaryInf$ !C, !D \fCenter !D \land !C $
\end{prooftree}
\end{ex}

\end{document}
